\section{連続的に変わる量の平均}\label{節:関数の平均}

\sectionref{節:定積分性質}で導入した関数の平均値を，ここでは温度のような連続的に変わる量へ使おう。まず，平均を積分で表す理由を振り返る。数 $u_1,\ldots,u_n$ の平均は和を個数で割ったものである。関数 $f$ の区間 $[a,b]$ における平均も，等間隔に値を取り，その平均の極限として考えられる。$\Delta x=\frac{b-a}{n}$ とすれば，
\[
\frac1n\sum_{k=1}^n f(a+k\Delta x)
 =\frac{1}{b-a}\sum_{k=1}^n f(a+k\Delta x)\Delta x.
\]
右辺の和はリーマン和なので，$f$ が連続なら極限は
\begin{equation}\label{式:関数の平均値}
\overline f=\frac{1}{b-a}\int_a^b f(x)\dd x
\end{equation}
となる。$f\geq0$ の場合，底辺 $b-a$，高さ $\overline f$ の長方形は，曲線の下と同じ面積になる。$f$ に負の値があっても，平均の定義は同じ式である。

\begin{example}{変化する温度の平均}
時刻 $0\leq t\leq12$ の温度を $T(t)=20+6\sin\frac{\pi t}{12}$ とする。この時間帯の平均温度を求めよ。
\end{example}

始めと終わりの温度はともに $20$ だが，途中の温度は $20$ より高い。両端の平均だけでは時間全体の平均にならない。定義に従って
\[
\begin{aligned}
\overline T
 &=\frac1{12}\int_0^{12}\left(20+6\sin\frac{\pi t}{12}\right)\dd t\\
 &=20+\frac12\left[-\frac{12}{\pi}\cos\frac{\pi t}{12}\right]_0^{12}
 =20+\frac{12}{\pi}.
\end{aligned}
\]
全時間を均等に数えて平均を取っていることが，$\frac1{12}$ に表れている。

連続な $f$ の最小値を $m$，最大値を $M$ とすると，$m\leq f(x)\leq M$ を積分して
\[
m\leq\overline f\leq M
\]
を得る。中間値の定理から，ある $c\in[a,b]$ で $f(c)=\overline f$ となる。これを積分の平均値の定理という。平均が実際のどこかの値と一致することは分かるが，一般にその点が区間の中央とは限らず，一つとは限らない。

均等に数えるのでない平均には，重みが必要になる。非負の連続関数 $w$ に対し $\int_a^b w(x)\dd x>0$ なら，重み付き平均は
\[
\frac{\int_a^b f(x)w(x)\dd x}{\int_a^b w(x)\dd x}
\]
である。分母は重みの総量であり，$w=1$ の場合が式\eqref{式:関数の平均値}になる。次に扱う重心は，位置の座標を面積や長さで重み付けして平均したものである。

